
\paragraph{\textbf{Theorem 4.11.}}
\SGARefTarget{sga3:xiv:4.11}
{\itshape Let \(G\) be a smooth algebraic group over an
algebraically closed field \(k\), and let \(\mathfrak g\) be its
Lie algebra. Then \(\mathfrak g\) is the union of the Lie algebras
\(\mathfrak b\) of the Borel subgroups \(B\) of \(G\).}

We may plainly suppose that \(G\) is connected. Let \(R\) be the
radical of \(G\), and put \(G'=G/R\). By \SGARefLink{sga3:xiv:4.9}{4.9}\textup{(b)} and \SGARefLink{sga3:xiv:4.10}{4.10},
it is enough to prove the theorem for \(G'\); thus we may suppose
that \(G\) is semisimple.

Let \(Z\) be the center of \(G\), which is its reductive center, and
put \(G'=G/Z\). The same argument, now using \SGARefLink{sga3:xiv:4.9}{4.9}\textup{(a)},
reduces the theorem to \(G'\). We may therefore suppose that \(G\)
is semisimple and adjoint.

Let \(B\) be a Borel subgroup of \(G\), let \(T\) be a maximal torus
of \(B\), hence of \(G\), and let \(\mathfrak b\) and
\(\mathfrak t\) be their Lie algebras. By \SGARefLink{sga3:xiv:3.18}{3.18}, \(T\) is a subgroup
of type \((\mathrm C)\) of \(G\), that is, \(\mathfrak t\) is a
Cartan subalgebra of \(\mathfrak g\). The union of the conjugates
of \(\mathfrak t\) is consequently dense in \(\mathfrak g\)
\textup{(\SGARefLink{sga3:xiii:theorem:5:1}{XIII~5.1}, \textup{(i)}\(\Rightarrow\)\textup{(vii)})}.
\emph{A fortiori}, the union of the conjugates of \(\mathfrak b\)
is dense in \(\mathfrak g\).

Let \(X\) be the closed subscheme of
\[
  G/B\times W(\mathfrak g)
\]
whose \(k\)-valued points are the pairs \((gB,x)\) such that
\[
  x\in\operatorname{Ad}(g)\mathfrak b
\]
\textup{(compare XIII~1)}. The morphism
\[
  \psi\colon X\longrightarrow W(\mathfrak g)
\]
induced by the second projection is proper because \(G/B\) is
proper over \(k\). We have just seen that it is dominant; it is
therefore surjective. This proves \SGARefLink{sga3:xiv:4.11}{4.11}.

The only result of this section that we shall use in the remainder
of this exposé is the following corollary.

\paragraph{\textbf{Corollary 4.12.}}
\SGARefTarget{sga3:xiv:4.12}
{\itshape Let \(k\) be an infinite field, let \(G\) be a smooth
algebraic group over \(k\), and let \(T\) be a maximal torus of
\(G\). Write \(\mathfrak g\) and \(\mathfrak t\) for their Lie
algebras. Let
\[
  u\colon G\longrightarrow\operatorname{GL}(V)
\]
be a linear representation, where \(V\) is finite-dimensional over
\(k\), and let
\[
  u'\colon\mathfrak g\longrightarrow\mathfrak{gl}(V)
\]
be the induced representation of Lie algebras. Then the minimum of
the nullity of \(u'(x)\), for \(x\in\mathfrak g\), is attained at
an element \(x\in\mathfrak t\).}

We reduce immediately to the case where \(k\) is algebraically
closed. We may plainly suppose that \(G\) is connected, and, after
dividing \(G\) by \((\ker u)_{\mathrm{red}}\), that \(G\) is
affine. Using \SGARefLink{sga3:xiv:4.11}{4.11} and the conjugacy theorem for the maximal tori
of \(G\), we reduce to the further case in which \(G\) is solvable.

Then \(G\) is a semidirect product \(T\cdot U\), where \(U\) is the
unipotent part of \(G\), a smooth connected unipotent group
\textup{\emph{(Bible)}, 6, Theorem~3}. Thus, as a vector space,
\[
  \mathfrak g=\mathfrak t\oplus\mathfrak u,
  \qquad
  \mathfrak u=\operatorname{Lie}(U).
\]
By the Lie--Kolchin theorem
\textup{\emph{(Bible)}, 6, \SGARefLink{sga3:xiv:7:lemma1}{Theorem~1}}, the representation space
\(V\) has a composition series by stable subspaces \(V_i\), with
one-dimensional quotients
\[
  W_i=V_i/V_{i+1}.
\]
For every \(i\) there is therefore an induced representation
\[
  u_i\colon G\longrightarrow
  \operatorname{GL}(W_i)=\mathbb G_m
\]
and a corresponding homomorphism of Lie algebras
\[
  u_i'\colon\mathfrak g\longrightarrow k.
\]
For every \(x\in\mathfrak g\), the nullity of \(u'(x)\) is the
number of indices \(i\) for which \(u_i'(x)=0\).

Since \(U\) is unipotent, every \(u_i\) is trivial on \(U\), and
hence every \(u_i'\) vanishes on \(\mathfrak u\). If
\[
  x=t+v,\qquad t\in\mathfrak t,\quad v\in\mathfrak u,
\]
then \(u_i'(x)=u_i'(t)\) for every \(i\). Thus \(u'(x)\) and
\(u'(t)\) have the same nullity, and \SGARefLink{sga3:xiv:4.12}{4.12} follows.
